\section{Sequential Multi-Agent Workflow (Architect / Implementer / Post-Implementation Review Agent)}

The workflow adopted in Part 2 separates planning, implementing, and reviewing into three distinct roles, coordinated through a strict sequence of hand-offs rather than concurrent execution, with a mandatory human gate between two of them. This is a different axis of change from Part 1, where a single agent carried out all three activities together, in one continuous exchange, rather than any one of them being delegated to a specialized role.

The Architect Agent goes first. Its job is to analyze a feature request against the existing codebase and produce a structured \textbf{\texttt{implementation\_plan.md}}, but it is explicitly barred from writing any implementation code itself --- its output is limited to architectural plans, interface definitions, and file-structure proposals. \Cref{lst:architect-agent-role} shows how this role's behavior is defined directly in \texttt{AGENTS.md}.

\begin{listing}[H]
\begin{center}
\begin{promptbox}
## 1. ROLE OF THE ARCHITECT AGENT

When acting as the Architect Agent, the agent MUST:
- Analyze the feature request against the existing codebase before proposing anything.
- Produce a structured implementation_plan.md artifact (see Section 11 for the required format).
- Include a mandatory Design Review section (Section 11.2) before listing any files or changes.
- NEVER write implementation code. Only produce architectural plans, interface definitions, and file-structure proposals.
- Flag any deviation from these rules explicitly in the plan and explain the justification.
- Identify which existing weaknesses (see Section 9) a new feature might worsen, and propose mitigations in the plan.

The Architect Agent hands off to the Implementer Agent only after the human has reviewed the Design Review section and explicitly approved the full plan.
\end{promptbox}
\end{center}
\caption{Excerpt of \texttt{AGENTS.md} defining the Architect Agent's role and behavioral constraints.}
\label{lst:architect-agent-role}
\end{listing}

Between the Architect and the Implementer sits a human review gate. The human is required to read the design review portion of the plan specifically; if the reasoning behind an architectural choice is vague, the human must ask for clarification, and if it is unclear who owns a given piece of state or which service is responsible for which logic, the plan must be rejected and revised rather than approved as written. Implementation cannot begin until this approval is given explicitly.

The Implementer Agent then executes, and its role is deliberately narrow: it does not design, decide, or deviate --- all of that is expected to have already happened in the plan. It is restricted to creating exactly the files listed under "files to create," modifying exactly the files listed under "files to modify," and implementing exactly the endpoints and components the plan describes, even where an unlisted addition might seem like a good idea on its own. If it encounters a situation during implementation that requires a design decision the plan never addressed, it is required to stop immediately, document what was expected against what was actually found, and wait for the Architect to produce an updated plan rather than resolving the gap itself. It is also bound by every rule in the charter regardless of whether the current plan happens to mention them, required to implement in a fixed dependency order --- backend models before services, services before routes, and so on through to frontend components --- and required to write the corresponding tests for anything it adds before considering the task finished. If it notices a bug or weakness while working that falls outside the current plan's scope, it is required to document it rather than fix it on the spot, since an unplanned fix is treated as its own task requiring its own plan.

Once the Implementer has finished and produced a \textbf{\texttt{walkthrough.md}}, the Post-Implementation Review Agent is activated. Its scope is deliberately narrow as well: it reads only the approved plan, the walkthrough, and the specific files the walkthrough lists as changed, and it does not re-check anything already enforced by the charter's rules, since the Implementer is already bound by those. Its purpose is to catch problems that could only emerge once code actually existed --- the same six-point checklist already described as part of the architectural conformance criterion, covering things like a component quietly taking on responsibilities beyond what was planned, unplanned coupling between modules meant to stay independent, and code paths left untested. It produces a \textbf{\texttt{review\_report.md}} rating each of the six concerns and closing with an overall verdict, after which a human decides whether to fix the issue immediately, defer it, or accept the feature as complete.

The first feature run through this full sequence was, fittingly, a fix to one of Part 1's own documented problems --- the \textbf{\texttt{MainPage.js}} monolith --- used as the initial test of whether the new process actually functioned as designed before being applied to newly added features.

